%%%%%%%%%%  scan idx 64  |  folio 51  %%%%%%%%%%
\textbf{Lemme 2.9.} \emph{Soit $\underline L=(L_i)$ un faisceau algébrique cohérent localement libre gradué, ayant un nombre fini de composantes non nulles. On a dans $\underline K(X)$ la formule~:}
\begin{equation}\tag{2.20}
 \underline E\left((\overset{p}{\otimes}\underline L)^{\mathrm{alt}}\right)=\wedge^p(\underline E(\underline L)).
\end{equation}
Dans cette formule, on considère $\overset{p}{\otimes}\underline L$ comme un faisceau où le groupe $\mathfrak S_p$ opère en tenant compte des graduations (qui introduisent donc des signes de façon bien connue, cf Cartan--Eilenberg)~; de plus, si $\underline M$ est un faisceau algébrique cohérent gradué sur $X$, on pose $\underline E(\underline M)=\sum_{i\geq0}(-1)^i\gamma(M_i)$.

Pour prouver la formule (2.20), on est ramené à prouver la formule analogue dans $\underline K(G)$, $G$ étant un groupe algébrique (un produit de groupes $\operatorname{Gl}(n,k)$ si l'on y tient) et $\underline L$ un $G$-module de dimension finie sur $k$ (cf. Chap.\ I, \S 2, exemple 2, pour la définition de $\underline K(G)$). Dans ce cas, cette formule a déjà été utilisée dans le théorème 1.4, pour montrer que le deuxième membre de (1.33) satisfait la même équation fonctionnelle que le premier membre~; la vérification en est élémentaire.
\begin{flushright}cqfd.\end{flushright}

Soit $f$ un morphisme de $X$ dans $Y$, $Y$ étant un espace algébrique. La notion d'image inverse de fibré vectoriel définit un homomorphisme $f^!:\underline K(\xi(Y))\longrightarrow\underline K(\xi(X))$ de $\lambda$-anneaux. Identifiant les fibrés vectoriels à des faisceaux algébriques localement libres, la notion d'image réciproque de fibré vectoriel correspond à la notion d'image réciproque de faisceau algébrique (déjà utilisée dans [3, n\textsuperscript{o} 8]). Dans le cas où $Y$ est non singulière, on définit plus généralement un homomorphisme de groupes~:
\begin{equation}\tag{2.21}
 f^!:\underline K(Y)\longrightarrow\underline K(X)
\end{equation}
